% OpenStax Calculus Volume 3 -- Section 3.3 Exercises: Arc Length and Curvature
% Exercises reproduced from OpenStax, licensed under CC BY 4.0.
% Source: https://openstax.org/books/calculus-volume-3/pages/3-3-arc-length-and-curvature
\documentclass{exercises}
\title{OpenStax Calculus Volume 3}
\subtitle{Section 3.3 Exercises: Arc Length and Curvature}
\sourceurl{https://openstax.org/books/calculus-volume-3/pages/3-3-arc-length-and-curvature}
\author{}
\date{}

\begin{document}
\maketitle


Find the arc length of the curve on the given interval.

\begin{enumerate}[start=1]
\item $\mathbf{r}(t)=t^{2}\mathbf{i}+14t\mathbf{j},\,0\le t\le7$. This portion of the graph is shown here: \\ \textit{[Figure: This figure is the graph of a curve beginning at the origin and increasing.]}
\item $\mathbf{r}(t)=t^{2}\mathbf{i}+(2t^{2}+1)\mathbf{j},\,1\le t\le3$
\item $\mathbf{r}(t)=\langle2\sin t,5t,2\cos t\rangle,0\le t\le \pi$. This portion of the graph is shown here: \\ \textit{[Figure: This figure is the graph of a curve in 3 dimensions. It is inside of a box. The box represents an octant. The curve begins in the upper right corner of the box and bends through the box to the other side.]}
\item $\mathbf{r}(t)=\langle t^{2}+1,4t^{3}+3\rangle,\,-1\le t\le0$
\item $\mathbf{r}(t)=\langle e^{-t}\cos t,e^{-t}\sin t\rangle$ over the interval $[0,\frac{\pi}{2}]$. Here is the portion of the graph on the indicated interval: \\ \textit{[Figure: This figure is the graph of a curve in the first quadrant. It begins approximately at $0.20$ on the $y$-axis and increases to approximately where $x=0.3$. Then the curve decreases, meeting the $x$-axis at $1.0$.]}
\item Find the length of one turn of the helix given by $\mathbf{r}(t)=\frac{1}{2}\cos t\mathbf{i}+\frac{1}{2}\sin t\mathbf{j}+\sqrt{\frac{3}{4}}\,t\,\mathbf{k}$.
\item Find the arc length of the vector-valued function $\mathbf{r}(t)=-t\mathbf{i}+4t\mathbf{j}+3t\mathbf{k}$ over $[0,1]$.
\item A particle travels once around a circle with the equation of motion $\mathbf{r}(t)=3\cos t\mathbf{i}+3\sin t\mathbf{j}+0\mathbf{k}$. Find the distance traveled around the circle by the particle.
\item Set up an integral to find the circumference of the ellipse with the equation $\mathbf{r}(t)=\cos t\mathbf{i}+2\sin t\mathbf{j}+0\mathbf{k}$.
\item Find the length of the curve $\mathbf{r}(t)=\langle \sqrt{2}t,e^{t},e^{-t}\rangle$ over the interval $0\le t\le1$. The graph is shown here: \\ \textit{[Figure: This figure is the graph of a curve in 3 dimensions. It is inside of a box. The box represents an octant. The curve begins in the upper left corner of the box and bends through the box to the bottom of the other side.]}
\item Find the length of the curve $\mathbf{r}(t)=\langle2\sin t,5t,2\cos t\rangle$ for $t\in[-10,10]$.
\item The position function for a particle is $\mathbf{r}(t)=a\cos(\omega t)\mathbf{i}+b\sin(\omega t)\mathbf{j}$. Find the unit tangent vector and the unit normal vector at $t=0$.
\item Given $\mathbf{r}(t)=a\cos(\omega t)\mathbf{i}+b\sin(\omega t)\mathbf{j}$, find the binormal vector $\mathbf{B}(0)$.
\item Given $\mathbf{r}(t)=\langle2e^{t},e^{t}\cos t,e^{t}\sin t\rangle$, determine the tangent vector $\mathbf{T}(t)$.
\item Given $\mathbf{r}(t)=\langle2e^{t},e^{t}\cos t,e^{t}\sin t\rangle$, determine the unit tangent vector $\mathbf{T}(t)$ evaluated at $t=0$.
\item Given $\mathbf{r}(t)=\langle2e^{t},e^{t}\cos t,e^{t}\sin t\rangle$, find the unit normal vector $\mathbf{N}(t)$ evaluated at $t=0$, $\mathbf{N}(0)$.
\item Given $\mathbf{r}(t)=\langle2e^{t},e^{t}\cos t,e^{t}\sin t\rangle$, find the unit binormal vector evaluated at $t=0$.
\item Given $\mathbf{r}(t)=t\mathbf{i}+t^{2}\mathbf{j}+t\mathbf{k}$, find the unit tangent vector $\mathbf{T}(t)$. The graph is shown here: \\ \textit{[Figure: This figure is the graph of a curve in 3 dimensions. It is inside of a box. The box represents an octant. The curve begins in the bottom left corner of the box and bends through the box to the upper left side.]}
\item Find the unit tangent vector $\mathbf{T}(t)$ and unit normal vector $\mathbf{N}(t)$ at $t=0$ for the plane curve $\mathbf{r}(t)=\langle t^{3}-4t,5t^{2}-2\rangle$. The graph is shown here: \\ \textit{[Figure: This figure is the graph of a curve above the $x$-axis. The curve decreases in the second quadrant, passes through the $y$-axis at $y=20$. Then it intersects the origin. The curve loops at the origin, increasing back through $y=20$ into the first quadrant.]}
\item Find the unit tangent vector $\mathbf{T}(t)$ for $\mathbf{r}(t)=3t\mathbf{i}+5t^{2}\mathbf{j}+2t\mathbf{k}$
\item Find the principal normal vector to the curve $\mathbf{r}(t)=\langle6\cos t,6\sin t\rangle$ at the point determined by $t=\pi/3$.
\item Find $\mathbf{T}(t)$ for the curve $\mathbf{r}(t)=(t^{3}-4t)\,\mathbf{i}+(5t^{2}-2)\,\mathbf{j}$.
\item Find $\mathbf{N}(t)$ for the curve $\mathbf{r}(t)=(t^{3}-4t)\,\mathbf{i}+(5t^{2}-2)\,\mathbf{j}$.
\item Find the unit normal vector $\mathbf{N}(t)$ for $\mathbf{r}(t)=\langle2\sin t,5t,2\cos t\rangle$.
\item Find the unit tangent vector $\mathbf{T}(t)$ for $\mathbf{r}(t)=\langle2\sin t,5t,2\cos t\rangle$.
\item Find the arc-length function $s(t)$ for the line segment given by $\mathbf{r}(t)=\langle3-3t,4t\rangle$. Write $\mathbf{r}$ as a parameter of $s$.
\item Parameterize the helix $\mathbf{r}(t)=\cos t\mathbf{i}+\sin t\mathbf{j}+t\mathbf{k}$ using the arc-length parameter $s$, from $t=0$.
\item Parameterize the curve using the arc-length parameter $s$, at the point at which $t=0$ for $\mathbf{r}(t)=e^{t}\sin t\mathbf{i}+e^{t}\cos t\mathbf{j}$.
\item Find the curvature of the curve $\mathbf{r}(t)=5\cos t\mathbf{i}+4\sin t\mathbf{j}$ at $t=\pi/3$. (Note: The graph is an ellipse.) \\ \textit{[Figure: This figure is the graph of an ellipse. The ellipse is oval along the $x$-axis. It is centered at the origin and intersects the $y$-axis at $-4$ and $4$.]}
\item Find the $x$-coordinate at which the curvature of the curve $y=1/x$ is a maximum value.
\item Find the curvature of the curve $\mathbf{r}(t)=5\cos t\mathbf{i}+5\sin t\mathbf{j}$. Does the curvature depend upon the parameter $t$?
\item Find the curvature $\kappa$ for the curve $y=x-\frac{1}{4}x^{2}$ at the point $x=2$.
\item Find the curvature $\kappa$ for the curve $y=\frac{1}{3}x^{3}$ at the point $x=1$.
\item Find the curvature $\kappa$ of the curve $\mathbf{r}(t)=t\mathbf{i}+6t^{2}\mathbf{j}+4t\,\mathbf{k}$. The graph is shown here: \\ \textit{[Figure: This figure is the graph of a curve in 3 dimensions. It is inside of a box. The box represents an octant. The curve has a parabolic shape in the middle of the box.]}
\item Find the curvature of $\mathbf{r}(t)=\langle2\sin t,5t,2\cos t\rangle$.
\item Find the curvature of $\mathbf{r}(t)=\sqrt{2}t\mathbf{i}+e^{t}\mathbf{j}+e^{-t}\mathbf{k}$ at point $P(0,1,1)$.
\item At what point does the curve $y=e^{x}$ have maximum curvature?
\item What happens to the curvature as $x\to \infty$ for the curve $y=e^{x}$?
\item Find the point of maximum curvature on the curve $y=\ln x$.
\item Find the equations of the normal plane and the osculating plane of the curve $\mathbf{r}(t)=\langle2\sin(3t),t,2\cos(3t)\rangle$ at point $(0,\pi,-2)$.
\item Find equations of the osculating circles of the ellipse $4y^{2}+9x^{2}=36$ at the points $(2,0)$ and $(0,3)$.
\item Find the equation for the osculating plane at point $t=\pi/4$ on the curve $\mathbf{r}(t)=\cos(2t)\mathbf{i}+\sin(2t)\mathbf{j}+t\mathbf{k}$.
\item Find the radius of curvature of $6y=x^{3}$ at the point $(2,\frac{4}{3})$.
\item Find the curvature at each point $(x,y)$ on the hyperbola $\mathbf{r}(t)=\langle a\cosh(t),b\sinh(t)\rangle$.
\item Calculate the curvature of the circular helix $\mathbf{r}(t)=r\sin(t)\mathbf{i}+r\cos(t)\mathbf{j}+t\,\mathbf{k}$.
\item Find the radius of curvature of $y=\ln(x+1)$ at point $(2,\ln 3)$.
\item Find the radius of curvature of the hyperbola $xy=1$ at point $(1,1)$.
\end{enumerate}

A particle moves along the plane curve $C$ described by $\mathbf{r}(t)=t\mathbf{i}+t^{2}\mathbf{j}$. Solve the following problems.

\begin{enumerate}[resume]
\item Find the length of the curve over the interval $[0,2]$.
\item Find the curvature of the plane curve at $t=0,1,2$.
\item Describe the curvature as $t$ increases from $t=0$ to $t=2$.
\end{enumerate}

The surface of a large cup is formed by revolving the graph of the function $y=0.25x^{1.6}$ from $x=0$ to $x=5$ about the $y$-axis (measured in centimeters).

\begin{enumerate}[resume]
\item \techrequired Use technology to graph the surface.
\item Find the curvature $\kappa$ of the generating curve as a function of $x$.
\item \techrequired Use technology to graph the curvature function.
\end{enumerate}

\end{document}
